% TOP_PROBLEM: {"id": 3191, "title": "Three-chromatic (0,2)-graphs", "category_id": 3, "category": {"id": 3, "name": "graph_theory", "display_name": "Graph Theory", "description": "Problems involving graphs, networks, and their properties.", "slug": "graph-theory", "order_index": 3, "created_at": "2026-07-31T15:26:25.671Z"}, "status": "open", "problem_number": "OPG-562", "set_id": 12}
% FIRST_POSTED: 2026-10-01
% ENTRY_KIND: statement_only
% UnsolvedMath ID 3191; source catalogue code OPG-562
% !TeX program = lualatex
% Definition and sources only; this file is not an LLM solution attempt.
\documentclass[11pt,a4paper]{article}
\usepackage[margin=25mm]{geometry}
\usepackage{fontspec}
\setmainfont{lmroman10-regular.otf}[BoldFont=lmroman10-bold.otf,
 ItalicFont=lmroman10-italic.otf,BoldItalicFont=lmroman10-bolditalic.otf]
\usepackage{amsmath,amssymb,mathtools}
\usepackage{unicode-math}
\setmathfont{latinmodern-math.otf}
\AtBeginDocument{\let\setminus\smallsetminus}
\usepackage{xurl,hyperref}
\hypersetup{colorlinks=true,urlcolor=blue}
\setcounter{secnumdepth}{0}
\setlength{\parindent}{0pt}
\setlength{\parskip}{5pt}
\setlength{\emergencystretch}{3em}
\begin{document}
\section*{Three-chromatic (0,2)-graphs}\label{3191}
{\small UnsolvedMath ID 3191 · OPG-562 · Graph Theory · OpenGarden}
\subsection{Definitions and mathematical statement}
Let $G=(V,E)$ be a finite simple undirected graph, and write $N(v)=\{u\in V:uv\in E\}$ for the open neighborhood of $v$. Impose the common-neighbor condition
\[
         |N(u)\cap N(v)|\in\{0,2\}
         \qquad\text{for all distinct }u,v\in V.
\]
The condition applies to adjacent and nonadjacent pairs alike. It does not require that every pair have two common neighbors, nor does it require regularity or a prescribed order.

The question is whether a graph satisfying this condition can have chromatic number exactly three. Explicitly, there must exist a map $c:V\to\{1,2,3\}$ with $c(u)\ne c(v)$ on every edge, while no map with just two colors has this property. Equivalently, the graph is three-colorable and contains an odd cycle. Neither merely having a displayed three-coloring nor using all three labels in a nonminimal coloring establishes this requirement.

The existence question concerns finite graphs; an infinite example would not answer it. Connectivity can be imposed without changing existence: a three-chromatic graph has a three-chromatic connected component, and each component inherits the common-neighbor condition. A positive answer should specify or establish such a finite graph and verify both requirements. A negative answer would show that every finite graph with the displayed common-neighbor property is either bipartite or requires at least four colors.
\subsection{Short English statement}
Is there a finite graph requiring exactly three colors in which every pair of distinct vertices has either zero or two common neighbors?
\subsection{Sources}
\begin{itemize}
\item[S1] ULAM.AI and the UnsolvedMath Contributors, supplied problems.json, record 3191 (OPG-562), OpenGarden. Catalogue identity and source transcription; CC BY 4.0.
\url{https://huggingface.co/datasets/ulamai/UnsolvedMath}.
\item[S2] Open Problem Garden, Three-chromatic (0,2)-graphs. Original problem and discussion.
\url{http://www.openproblemgarden.org/op/three_chromatic_0_2_graphs}.
\end{itemize}
\end{document}
